\section{Centering and deleting a small edge mass}\label{q:analytic}

We first compare the progression correlation with a sum whose prime
factors have mean zero in the independent residue model.  The short
Fourier estimate controls the terms introduced by this centering.
We then show that restrictions on padding mass and prime degrees, together
with any suitably rare residue event, remove only a small total mass.
These two estimates prepare the correlation for the graph argument.

\subsection{Centering the prime divisibility conditions}

We first bound the terms other than full divisibility in the expansion
of $C_j$ from \eqref{q:centered-average}. Each such term leaves a
nonempty product of primes to be averaged as a shift. Our target is
an error that remains small after summing over all $O(L/\eta)$ bins.

\begin{lemma}[Centering estimate]\label{q:center}
With the parameters and prime supplies fixed above, for every sufficiently
large real $X$ one has
\begin{equation}\label{q:center-estimate}
 \left|\sum_j C_j-S_0F(X)\right|
 \ll_{h,l}SV_*\left(
   \eta+X^{-1}+\eta^{-1}2^JL^{-1/20}\right).
\end{equation}
\end{lemma}

We prove this estimate after two Fourier bounds. The first concerns
rough integers; the second combines it with the short exponential-sum
theorem of Matom\"aki, Radziwi\l\l, and Tao. The resulting shift
average saves more than $L^{-1}$, which pays for the number of bins.
The fourth-moment treatment of the divisor exponential sum follows
\cite[Appendix~C]{Pilatte26}; we include the rough-number estimates and
their finite counting errors. Write $\mathrm e(t)=\exp(2\pi i t)$.

\begin{lemma}[Fourier bounds for rough integers]\label{q:rough-fourier}
Let $L$ be sufficiently large, let $h\ge1$ be an integer, and suppose
\[
 \tfrac12\exp(L^{0.995})\le D\le\exp(2L).
\]
Let $\mathcal Z\subset[D,2D)$ consist of integers with no prime factor
at most $\exp(L^{0.99})$. For
\[
 B(\theta)=\sum_{z\in\mathcal Z}\frac{\mathrm e(hz\theta)}z
 \qquad(\theta\in\mathbb R/\mathbb Z)
\]
one has, with absolute implied constants,
\begin{equation}\label{q:rough-bounds}
 \|B\|_\infty\ll L^{-0.99},\qquad
 \int_0^1|B(\theta)|^4\,d\theta\ll D^{-1}L^{-3.96}.
\end{equation}
\end{lemma}

\begin{proof}
Put $y=\exp(L^{0.99})$. We need upper bounds for the number of integers
in $[D,2D)$ avoiding every prime at most $y$, and for the number of triples
in $[D,2D)^3$ for which all four forms
\[
 x_1,\quad x_2,\quad x_3,\quad x_1+x_2-x_3
\]
avoid those primes. The last form is positive throughout this integer box.
In the first case put $m=1$, and in the second put $m=4$. For each prime $p$,
the proportion $\nu_m(p)$ of residue vectors where at least one form
vanishes satisfies
\begin{equation}\label{q:sieve-local-density}
 \nu_m(p)=\frac mp+O(p^{-2}).
\end{equation}
For the four forms, every pair of defining hyperplanes is independent over
every prime field, including the field with two elements. Inclusion and
exclusion of pairs proves \eqref{q:sieve-local-density}. Moreover
$\nu_m(p)<1$ for every $p$: the residue $1$ in the first case and the vector
$(1,1,1)$ in the second avoid all the forms.

For completeness, a finite inclusion--exclusion argument suffices here.
Let $r_0=2\lceil500\log L\rceil$ and truncate inclusion--exclusion of the
bad-prime events at this even order. This gives an upper bound for the
proportion avoiding all the events. For any selected prime set of product
$r$, the Chinese remainder theorem gives residue density
$\prod_{p\mid r}\nu_m(p)$. Its normalized count in the interval or box
differs from this density by $O(r/D)$ when $r\le D$. Indeed, each residue
coordinate occurs $D/r+O(1)$ times; summing the product of these counts over
at most $r$ or $r^3$ residue vectors proves the assertion.
There are at most
$(r_0+1)(1+\pi(y))^{r_0}$ terms, each with $r\le y^{r_0}$.
The total boundary error is therefore at most
\begin{equation}\label{q:sieve-boundary-error}
 \frac1D\exp\bigl(O(L^{0.99}\log L)\bigr),
\end{equation}
which is smaller than every fixed negative power of $L$, by the lower bound
on $D$.

In the independent residue model the difference between the even truncated
sum and the full product is at most the next elementary symmetric sum.
Partial summation of the prime number theorem and
\eqref{q:sieve-local-density} give
\[
 \sum_{p\le y}\nu_m(p)\le(4+o(1))\log L.
\]
Thus that difference is at most
\[
 \frac{\bigl(\sum_{p\le y}\nu_m(p)\bigr)^{r_0+1}}{(r_0+1)!}
 \le
 \left(\frac{e\sum_{p\le y}\nu_m(p)}{r_0+1}\right)^{r_0+1}
 \ll L^{-100}.
\]
The full product satisfies
\[
 \prod_{p\le y}(1-\nu_m(p))\ll(\log y)^{-m}=L^{-0.99m};
\]
the finitely many small primes contribute a positive constant, and the
remaining factors follow by taking logarithms in
\eqref{q:sieve-local-density}. Together with
\eqref{q:sieve-boundary-error}, this proves the respective counts
$O(DL^{-0.99})$ and $O(D^3L^{-3.96})$.

The first count bounds $\sum_{z\in\mathcal Z}1/z$ and hence
$\|B\|_\infty$. Orthogonality expresses the fourth moment as the sum of
$(z_1z_2z_3z_4)^{-1}$ over
$z_1+z_2=z_3+z_4$ with all four variables in $\mathcal Z$.
The factor $h$ does not change this equality. Each triple determines $z_4$,
and each weight is at most $D^{-4}$. The second count proves the fourth
moment bound.
\end{proof}

The analytic input keeps the frequency fixed while averaging the starting
point of a short interval.

\begin{theorem}[Matom\"aki--Radziwi\l\l--Tao,
\cite{MRT15}, Theorem~1.3]\label{q:MRT-short}
For every pair of real numbers $10\le H\le Z$,
\[
 \sup_{\alpha\in\mathbb R}
 \int_0^Z
 \left|\sum_{x\le n\le x+H}\lambda(n)\mathrm e(\alpha n)\right|\,dx
 \ll HZ\left(\frac{\log\log H}{\log H}
                  +(\log Z)^{-1/700}\right),
\]
with an absolute implied constant. The supremum is outside the integral.
\end{theorem}

We next combine this theorem with Lemma~\ref{q:rough-fourier}. The conclusion
is uniform over arbitrary subsets of the rough integers, so it permits all
the bin and prime-supply restrictions imposed in
\eqref{q:centered-average}.

\begin{lemma}[Correlation averaged over rough shifts]\label{q:rough-shifts}
Fix integers $h,l\ge1$ and $b_0\in\mathbb Z$. Let $A\ge1000$, let $L$ be
sufficiently large, and suppose $Y\ge\exp(L^A/2)$. Let $D$ be a positive
integer and let $\mathcal Z\subset[D,2D)$ satisfy the assumptions of
Lemma~\ref{q:rough-fourier}. Then
\begin{equation}\label{q:rough-correlation}
 \left|\frac1Y\sum_{1\le v\le Y}
       \sum_{z\in\mathcal Z}\frac1z
       \1_{v\equiv b_0\pmod l}\lambda(v)\lambda(v+hz)\right|
 \ll_{h,l}L^{-21/20}.
\end{equation}
\end{lemma}

\begin{proof}
For each integer $1\le v\le Y$ define
\begin{align*}
 F_v(\theta)&=\sum_{m=1}^D
  \lambda(v+m)\1_{v+m\equiv b_0\pmod l}\mathrm e(\theta m),\\
 G_v(\theta)&=\sum_{a=1}^{(2h+1)D}
  \lambda(v+a)\mathrm e(-\theta a).
\end{align*}
We first claim that, uniformly in $\theta$,
\begin{equation}\label{q:progression-short-sum}
 \sum_{1\le v\le Y}|F_v(\theta)|\ll YD L^{-0.8}.
\end{equation}
For integer $D$, every $x\in(v,v+1)$ has precisely
$v+1,\ldots,v+D$ as the integers in $[x,x+D]$.
Apply Theorem~\ref{q:MRT-short} with $H=D$ and
$Z=\lfloor Y\rfloor+1$, discarding the unused initial unit interval.
The restriction to the residue class follows from
\[
 \1_{n\equiv b_0\pmod l}
   =\frac1l\sum_{a=0}^{l-1}\mathrm e\left(\frac{a(n-b_0)}l\right).
\]
Its coefficients have total absolute value one. After the unimodular
factor arising from $n=v+m$ is removed, this applies the same theorem to
the frequencies $\theta+a/l$. In particular, no assumption that $b_0$ is
coprime to $l$ is needed. The scale hypotheses imply $10\le D\le Z$ and
\[
 \frac{\log\log D}{\log D}+(\log Z)^{-1/700}
 \ll L^{-0.8},
\]
which proves \eqref{q:progression-short-sum}.

Average the sum on the left of \eqref{q:rough-correlation}, before taking
its absolute value, over translations by $m=1,\ldots,D$. Each translation
changes at most $2m$ prefix terms for each $z$. By
Lemma~\ref{q:rough-fourier}, the normalized endpoint error is at most
\[
 O\left(\frac D Y\sum_{z\in\mathcal Z}\frac1z\right)=O(D/Y).
\]
Let $B$ be the polynomial in Lemma~\ref{q:rough-fourier}. Orthogonality now
expresses the translated average as
\begin{equation}\label{q:short-convolution}
 \frac1{YD}\sum_{1\le v\le Y}
   \int_0^1B(\theta)F_v(\theta)G_v(\theta)\,d\theta.
\end{equation}
Indeed, expanding the integral imposes $a=m+hz$. Since
$1\le m\le D$ and $D\le z<2D$, this value of $a$ lies in the full range
used to define $G_v$.

On the set where $|B|\le L^{-1.05}$, Parseval and Cauchy--Schwarz give
\[
 \int_0^1|F_vG_v|\le\|F_v\|_2\|G_v\|_2
 \le\sqrt{D}\sqrt{(2h+1)D}\ll_h D.
\]
Its contribution to \eqref{q:short-convolution} is therefore
$O_h(L^{-1.05})$. The complementary set, denoted by $\mathcal E$, satisfies
\[
 |\mathcal E|\le L^{4.2}\int_0^1|B|^4
 \ll D^{-1}L^{0.24}.
\]
On this set use $\|B\|_\infty\ll L^{-0.99}$,
$|G_v|\le(2h+1)D$, and \eqref{q:progression-short-sum}. Its contribution is
at most
\[
 \frac1{YD}|\mathcal E|\,L^{-0.99}(2h+1)D\,(YDL^{-0.8})
 \ll_h L^{0.24-0.99-0.8}=L^{-1.55}.
\]
Finally $D/Y\le\exp(2L-L^A/2)$ is negligible. This proves
\eqref{q:rough-correlation}.
\end{proof}

\begin{proof}[Proof of \Cref{q:center}]
The full-divisibility contribution is \eqref{q:full-divisibility}.
For every other term in the expansion of \eqref{q:centered-average},
let $I\subset\{1,\ldots,J\}$ be the nonempty set of
prime supplies whose factors have been replaced by $-1/p$. Put
\[
 z=\prod_{i\in I}p_i,\qquad d'=d/z.
\]
For fixed $q,d',j$, substitute $n=qd'v$ and put $Y=T_j/(qd')$.
After extracting the sign $(-1)^{|I|}$ and the factor $u(q)/(qd')$, the
remaining expression is
\begin{equation}\label{q:partial-centering}
 \frac1Y\sum_{1\le v\le Y}\sum_z\frac1z
    \1_{v\equiv bz\pmod l}\lambda(v)\lambda(v+hz),
\end{equation}
where $z$ ranges over the products from the supplies indexed by $I$ that
also satisfy the bin condition. This substitution is valid in every
residue class, since $qd'$ is invertible modulo $l$.

The allowed $z$ lie in an interval of logarithmic width $\eta<\log2$.
Split them by residue modulo $l$ and by the at most two dyadic intervals
$[D,2D)$ that meet this interval. In each resulting piece the residue
$bz\pmod l$ is constant. Unique prime factorization and the disjointness
of the supplies show that each integer $z$ occurs at most once. Since
$I$ is nonempty, all its prime factors are at least
$H_0=\exp(L^{0.995})$, and
\[
 D\ge H_0/2,\qquad D\le\exp(2L).
\]
Furthermore, if the piece is nonempty, its bin condition implies
$Y>Xe^{-\eta}z$, so $\log Y\ge L^A/2$ for large $L$.
Lemma~\ref{q:rough-shifts} therefore bounds
\eqref{q:partial-centering} by $O_{h,l}(L^{-21/20})$.

For each fixed $I$ and bin, the sum of the extraction weights is at most
\[
 \sum_q\frac{u(q)}q\sum_{d'}\frac1{d'}
 \le S\prod_{i\notin I}V_i\le SV_*,
\]
where the last inequality uses $V_i>1$. There are
$O(L/\eta)$ bins and fewer than $2^J$ nonempty sets $I$. Hence the total
contribution of the partial centering terms is
\[
 O_{h,l}\bigl(SV_*\,2^J\eta^{-1}L^{-1/20}\bigr).
\]
Together with the full divisibility term this proves
\eqref{q:center-estimate}.
\end{proof}

\subsection{The distribution of padding divisors}

The centered averages still have occasional sites at which many padding
choices contribute to one bin.  We next prove that deleting such sites
has small total cost.  For a site $n$, a label $d\in\mathcal D$, and one of
the bin indices $j$, define
\begin{equation}\label{q:rho}
 \rho_j(n,d)=\frac1{w(n)}
   \sum_{\substack{q\in\mathcal Q,\ q\mid n\\
          j\eta\le\log(qd)<(j+1)\eta}}u(q).
\end{equation}
Here the sum includes all squarefree padding divisors in the bin, without
the restriction on $\omega(q)$ imposed in $\mathcal S_j$.
The denominator is positive because the divisor $1$ is always present.
Thus $0\le\rho_j\le1$ and $\sum_j\rho_j\le1$.
All these definitions also make sense in the auxiliary residue model.
At a translated site $n+a$, divisibility by $p$ means
$R_p+a\equiv0\pmod p$, using the same residue coordinate at every site.

Define the tilted expectation at a single site by
\[
 \E^* H=\frac{\E(w(n)H)}S.
\]
This is a probability expectation since
\[
 \E w(n)=\prod_{r\in Q}(1+4/r)=S.
\]
The tilt preserves independence of the $Q$ coordinates and changes the
probability of $r\mid n$ from $1/r$ to $5/(r+4)$.

\begin{lemma}[Padding anti-concentration]\label{q:padding}
For every $d\in\mathcal D$ and the bin width $\eta=e^{-J}$,
\begin{equation}\label{q:padding-square}
 \E^*\sum_j\rho_j(n,d)^2\ll_{h,l}L^{-1}.
\end{equation}
The constant is independent of $d$, $W$, and the location of the bins.
\end{lemma}

\begin{proof}
Conditionally on the primes of $Q$ dividing $n$, choose $q_1,q_2$
independently with probabilities
\[
 \Pbb(q_a=q\mid n)=\frac{u(q)\1_{q\mid n}}{w(n)}
 \qquad(a=1,2).
\]
Each available prime is selected with probability $4/5$.  Under the
combined tilted law let $Z=\log(q_1/q_2)$.  A prime $r$ contributes
$+\log r$ or $-\log r$ with probability $4/[5(r+4)]$ each, and contributes
zero otherwise.  The contributions are independent, so its characteristic
function, with angular frequency $t$, is
\begin{equation}\label{q:padding-characteristic}
 \phi(t)=\E^*e^{itZ}
 =\prod_{r\in Q}\left(1-\frac8{5(r+4)}
                              (1-\cos(t\log r))\right).
\end{equation}
Every factor is nonnegative: even at $r=2$ it is at least $7/15$.
The inequality $1-x\le e^{-x}$ therefore gives
\[
 0\le\phi(t)\le
 \exp\left(-\frac85\sum_{r\in Q}
                     \frac{1-\cos(t\log r)}{r+4}\right).
\]
Replacing $r+4$ by $r$ changes the exponent by $O(1)$, uniformly in $t$.
For $|t|\le2$, partial summation of the fixed-modulus prime number theorem
in \eqref{q:pnt} gives
\begin{align}
 \sum_{r\in Q}\frac{1-\cos(t\log r)}r
 &=\frac34\int_{\log2}^L\frac{1-\cos(ty)}y\,dy+O_{h,l}(1)
       \notag\\
 &=\frac34\log(1+L|t|)+O_{h,l}(1).\label{q:padding-pnt}
\end{align}
For completeness, in the partial summation the derivative of
$(1-\cos(t\log x))/(x\log x)$ is $O(1/(x^2\log x))$ uniformly for
$|t|\le2$.  Its product with the prime-number-theorem error is integrable
on $[2,\infty)$.  Deleting the finitely many prime divisors of $lh$ also
costs $O_{h,l}(1)$.  To verify the second equality, set $u=|t|y$.
The integral over $u\le1$ is bounded, and integration by parts bounds
$\int_1^U\cos u\,du/u$ uniformly in $U\ge1$.  The lower endpoint
$|t|\log2$ stays bounded.  These observations also cover $t=0$.
It follows that
\begin{equation}\label{q:padding-decay}
 0\le\phi(t)\ll_{h,l}(1+L|t|)^{-6/5}\qquad(|t|\le2).
\end{equation}

Conditionally on $n$, the probability that both divisors lie in the
$j$th bin is $\rho_j(n,d)^2$. Divisors in the same bin have
$|Z|\le\eta\le1$. Consequently the left side of
\eqref{q:padding-square} is at most
$\Pbb^*(|Z|\le1)$.  The function
$\kappa(x)=(\sin x/x)^2$, with $\kappa(0)=1$, satisfies
\[
 \1_{[-1,1]}(x)\le\frac{\kappa(x)}{\sin^2 1},\qquad
 \kappa(x)=\frac14\int_{-2}^2(2-|t|)e^{itx}\,dt.
\]
Taking expectations, using \eqref{q:padding-decay}, and integrating gives
\[
 \Pbb^*(|Z|\le1)
 \ll_{h,l}\int_{-2}^2(1+L|t|)^{-6/5}\,dt
 \ll_{h,l}L^{-1}.
\]
The common translation $\log d$ of the two padding logarithms disappeared
from $Z$, which proves the asserted uniformity.
\end{proof}

\subsection{Deleting large degrees and prescribed rare sites}

We formulate the deletion step with an explicit input for a family of rare
sites.  The following section will construct that family from short
paths.  For each bin $j$, choose a Boolean predicate of the residue coordinates
and write $n\in\mathcal B_j$ when it holds at site $n$.  Translation
covariance means that $n+a\in\mathcal B_j$ is obtained from the same
predicate by adding $a$ to every residue coordinate.  We assume
\begin{equation}\label{q:rare-input}
 \Pbb(n\in\mathcal B_j)\le
           \exp\bigl(-\tfrac12L^{1-\delta}\bigr),
\end{equation}
and that the event is represented by an AND--OR--NOT circuit of depth at
most $6$ and size at most $\exp(L^3)$, whose inputs are equalities in the
prime residue coordinates and the coordinate modulo $l$.
The event uses no value of the Liouville function.  Define
$\sigma_j(n)=\1_{n\notin\mathcal B_j}$.  For any prime set $R$,
write $\omega_R(n)=\#\{p\in R:p\mid n\}$.  Define
\begin{align}
 \tau_j(n,d)&=
 \1_{\omega_Q(n)\le400\log L}\,
 \1_{\rho_j(n,d)\le K/L}\,\sigma_j(n),\notag\\
 \chi_j(n,d)&=
 \tau_j(n,d)\,\1_{\omega_P(n)\le6WJ}.
 \label{q:tau}
\end{align}
The function $\tau_j$ omits the $P$-degree cut: the graph argument will
integrate the centered $P$ coordinates before applying that restriction
as a projection. At present we impose all four cuts at
both ends of each summand, defining
\begin{align}
 C_j^\circ={}&\frac1{T_j}\sum_{1\le n\le T_j}
 \sum_{(d,q)\in\mathcal S_j}
 u(q)\1_{q\mid n}\1_{n\equiv bqd\pmod l}
 \prod_{p\mid d}\left(\1_{p\mid n}-\frac1p\right)
 \lambda(n)\lambda(n+hqd)\notag\\[-2pt]
 &\hspace{35mm}\cdot\chi_j(n,d)\chi_j(n+hqd,d).
 \label{q:cut-averages}
\end{align}

\begin{lemma}[Deletion cost]\label{q:cuts}
Suppose the family $\mathcal B_j$ satisfies the covariance, probability,
and circuit hypotheses just stated.  Then
\begin{equation}\label{q:cuts-bound}
 \frac1{SV_*}\sum_j|C_j-C_j^\circ|
 \ll_{h,l}
 2^J\bigl(K^{-1}+L^{-100}+e^{-2WJ}\bigr)
       +\exp(-L^{0.9}).
\end{equation}
This estimate holds for the integer averages defining $C_j$ and
$C_j^\circ$, for every sufficiently large real $X$.
\end{lemma}

\begin{proof}
We first bound the deletion costs in the auxiliary product model, after
dropping the Liouville factors and the progression indicator, and replacing
the absolute centered product by
\[
 a_d^+(n)=\prod_{p\mid d}\left(\1_{p\mid n}+\frac1p\right).
\]
At the end of the proof we transfer precisely these nonnegative costs
to integer averages by Lemma~\ref{q:finite-law}.  No assertion about
untruncated integer moments of $w$ is needed.

\paragraph{Either endpoint.}
For a fixed numerical pair $(d,q)$, translation by $hqd$ preserves the
auxiliary product law.  Moreover divisibility by $q$ and by each prime
of $d$ is unchanged by this translation.  Thus an upper-endpoint failure
has the same bound as a lower-endpoint failure.  This argument does not
assert that the full weight $w$ is constant across an edge.  It suffices
to analyze one endpoint and multiply the resulting bound by $2$.

\paragraph{The two padding cuts.}
The $P$ coordinates are independent of the $Q$ coordinates and
$\E a_d^+=2^J/d$.  Enlarge the eligible padding divisors to all divisors
in their respective bins.  The total costs of the first two cuts, at
one endpoint and summed over bins, are at most
\begin{equation}\label{q:padding-deletion}
 2^JS\sum_{d\in\mathcal D}\frac1d
 \left\{\Pbb^*(\omega_Q>400\log L)
       +\E^*\sum_j\rho_j(n,d)
                    \1_{\rho_j(n,d)>K/L}\right\}.
\end{equation}
Indeed the sum of padding weights in a bin is $w(n)\rho_j(n,d)$,
and $\sum_j\rho_j\le1$.  Under the tilted law the mean of $\omega_Q$
is at most $4\log L$ for large $L$.  Independence and exponential Markov
therefore give
\[
 \Pbb^*(\omega_Q>400\log L)
 \le\exp\bigl((-400+4(e-1))\log L\bigr)\ll L^{-100}.
\]
Also $x\1_{x>K/L}\le(L/K)x^2$ for $x\ge0$.  Applying
Lemma~\ref{q:padding} bounds the second term in braces by
$O_{h,l}(K^{-1})$.  Since $\sum_d1/d=V_*$, this gives the first two
terms of \eqref{q:cuts-bound}.

\paragraph{The $P$-degree cut.}
For fixed $d$, tilt the $P$ coordinates by $a_d^+$ divided by its mean
$2^J/d$.  Primes outside $d$ retain their independent Bernoulli laws,
whereas the selected $J$ primes contribute at most $J$ to the count.
As $\sum_{p\in P}1/p=\sum_iV_i\le2WJ$, exponential Markov gives
\[
 \Pbb_{d,+}(\omega_P>6WJ)
 \le\exp\bigl(-6WJ+J+2(e-1)WJ\bigr)
 \le e^{-2WJ}\qquad(W\ge10).
\]
The padding coordinates remain independent.  Summing their weights
$u(q)/q$ over all eligible pairs and bins is at most $S$ for each $d$.
Thus this deletion costs at most $2^JSV_*e^{-2WJ}$ per endpoint.

\paragraph{The rare sites.}
For one bin, the entire absolute row weight is bounded by
\[
 B(n)=w(n)\prod_{i=1}^J(\omega_{P_i}(n)+V_i).
\]
To see this, sum all padding divisors for each $d$, and then sum each
prime slot separately.  Independence gives
\[
 \E B(n)^2
 \le\prod_{r\in Q}(1+24/r)
       \prod_{i=1}^J(4V_i^2+V_i)
 \le C_{h,l}L^{C}(18W^2)^J
 \le C_{h,l}L^{C'}.
\]
Here $C,C'$ are absolute: the first product is polynomial in $L$ by
partial summation, and
$J\log(18W^2)\le\delta\log L\,\log(18W^2)/(6W)$ is bounded by an
absolute multiple of $\log L$ for $W\ge10$.
Cauchy--Schwarz and \eqref{q:rare-input} therefore bound a single-bin
rare-site cost by
\[
 C_{h,l}L^{C'/2}\exp\bigl(-\tfrac14L^{1-\delta}\bigr).
\]
There are $O(L/\eta)$ bins, and $SV_*\ge1$.
Their total contribution is at most $\exp(-L^{0.9})$ for sufficiently
large $L$.

\paragraph{Transfer to integer averages.}
We now justify applying Lemma~\ref{q:finite-law} to the preceding costs.
This also ensures that no independence claim is being made for arbitrary
integer functions.  Use a union bound for the two endpoints, and for
the following four failures at either endpoint:
\[
 \omega_Q>400\log L;\qquad
 \omega_Q\le400\log L\text{ and }\rho_j>K/L;\qquad
 n\in\mathcal B_j;\qquad \omega_P>6WJ.
\]
This is the same deletion event, with the padding-mass failure tested
only after the low-$Q$ cut passes.  In the product model its cost is
bounded by the estimates already proved, since dropping the extra
low-$Q$ condition can only increase a nonnegative cost.

Fix a numerical $j,d,q$ and expand $a_d^+$ as a sum of divisibility
indicators.  There are $2^J$ terms, with coefficients at most $1$.
The factor $\1_{q\mid n}$ is a conjunction of its prime divisibilities.
Each degree failure is an OR over subsets of primes of the required
cardinality, testing that all their residues vanish.  Since there are
at most $e^L$ available primes, the low-$Q$ degree test has size
$\exp(O(L\log L))$.  The $P$ threshold has the same bound, using
$6WJ\le\delta\log L$.

For the padding-mass failure with low $Q$ degree, list all possible exact
sets of at most $400\log L$ primes of $Q$ dividing the site.  For each
set, a conjunction asserts both the indicated presences and all other
absences.  On that state, both $w$ and $\rho_j$ are fixed numbers; retain
the state precisely when $\rho_j>K/L$.  There are
$\exp(O(L\log L))$ such states, each specified by at most $e^L$ tests.
The remaining event $\mathcal B_j$ has the circuit representation assumed
in the statement.  Intersecting any of these events with the indicated
prime divisibilities has depth at most $10$ and size at most
$\exp(L^4)$ for large $L$.  Translation to an upper endpoint changes
only the tested residues, not these bounds.

Finally, the numerical pairs satisfy $dq\le\exp(100L+\eta)$.
The disjoint prime supports of $d$ and $q$ determine both from their
product, so there are $\exp(O(L))$ eligible pairs in all bins.  Also
$u(q)\le4^{100\log L}$, and the expansion of $a_d^+$ has total
coefficient mass at most $2^J$.  The total absolute coefficients in all
these comparisons are therefore at most $\exp(O(L^4))$; even the looser
bound permits all listed state expansions.  The averaging interval has
length $\lfloor T_j\rfloor\ge\exp(\tfrac12L^A)$.
Lemma~\ref{q:finite-law} makes the aggregate comparison error at most
$\exp(-L^9)$.  Replacing the integer-length normalization by $T_j$ costs
at most $T_j^{-1}\exp(O(L^4))$, which is smaller still.
Together with the product estimates, these errors prove
\eqref{q:cuts-bound}.
\end{proof}
